\section{Design of Long-Run Workflows}
\subsection{Loop}
\begin{frame}{Loop Roadmap}
\begin{itemize}
\item \textbf{Loop}
\begin{itemize}
\item Act, observe and decide
\item Revise and critique
\item Carry findings and verify repairs
\item Accept, stop and recover
\end{itemize}
\end{itemize}
\end{frame}
\begin{frame}{Act, Observe, Decide}
\begin{center}\resizebox{.91\textwidth}{!}{\begin{tikzpicture}[b/.style={draw=blue,fill=blue!5,rounded corners,text width=2.9cm,minimum height=.8cm,align=center,font=\small},a/.style={-{Stealth},thick}]
\node[b](p)at(0,1.2){Plan next step};
\node[b](t)at(4,1.2){Allowed tool call};
\node[b](o)at(8,1.2){Observe + decide};
\node[b](r)at(2,-0.8){Continue within budget};
\node[b](e)at(8,-0.8){Return or pause};
\draw[a](p)--(t);\draw[a](t)--(o);\draw[a](o.south)--(r.east);\draw[a](r.west)-| (p.south);\draw[a](o)--(e);\end{tikzpicture}
}\end{center}
\begin{itemize}
\item A failed tool call leaves work to do.
\item Use the observation to choose the next action.
\item Retry temporary errors within a fixed limit; invalid evidence must not advance the case.
\end{itemize}
\end{frame}
\begin{frame}{Revision Loop}
\begin{center}\resizebox{.96\textwidth}{!}{\begin{tikzpicture}[
  box/.style={draw=blue,fill=blue!5,rounded corners=2pt,text width=2.65cm,minimum height=1.4cm,align=center,font=\footnotesize,inner sep=4pt},
  flow/.style={-{Stealth[length=4pt]},line width=.7pt,draw=black!60},
  lab/.style={font=\scriptsize,text=black!60,fill=white,inner sep=2pt}]
\node[box] (g) at (1.6,0) {\textbf{Generator}\\candidate + tests};
\node[box,draw=forestgreen,fill=forestgreen!5] (t) at (5.05,0) {\textbf{Code gates}\\required checks};
\node[box,draw=darkred,fill=red!4] (c) at (8.5,0) {\textbf{Fresh Critic}\\inspect evidence};
\node[box,draw=forestgreen,fill=forestgreen!5] (o) at (11.95,0) {\textbf{Orchestrator}\\revise, accept\\or stop};
\draw[flow] (g)--(t);\draw[flow] (t)--(c);\draw[flow] (c)--(o);
\draw[flow] (o.south)--(11.95,-1.35)--node[below,lab]{repair packet: findings + required evidence}(1.6,-1.35)--(g.south);
\node[box,text width=3.3cm,minimum height=1.0cm] (claim) at (2.1,-3.0) {\textbf{Generator's claim}\\\lq\lq Costs included\rq\rq};
\node[box,text width=3.3cm,minimum height=1.0cm] (inspect) at (6.8,-3.0) {\textbf{Critic's check}\\saved trades + backtest};
\node[box,text width=3.3cm,minimum height=1.0cm,draw=darkred,fill=red!4] (finding) at (11.5,-3.0) {\textbf{Finding}\\costs were not deducted};
\draw[flow] (claim)--(inspect);\draw[flow] (inspect)--(finding);
\node[font=\scriptsize,text=black!55] at (6.8,-3.85) {Illustrative review};
\end{tikzpicture}
}\end{center}
\begin{itemize}
\item Give the Critic saved trades and backtest outputs to inspect.
\item A fresh Critic each round can reduce anchoring on earlier responses.
\item Keep the mandate and acceptance rules fixed during revision.
\end{itemize}
\end{frame}
\begin{frame}{Findings and Repairs\handsup}
\begin{center}\resizebox{.94\textwidth}{!}{\begin{tikzpicture}[
  card/.style={draw=blue,fill=blue!4,rounded corners=3pt,align=left,font=\small,inner sep=7pt}]
\node[card] at (3.1,0) {
  \parbox[t][3.25cm][t]{5.65cm}{\raggedright
  \textbf{Timing Problem}\par\smallskip
  The signal uses day-$t$ closing prices.\par\smallskip
  The backtest trades at that same close.\par\smallskip
  The trade must follow signal availability.
  }
};
\node[card,draw=darkred,fill=red!3] at (9.9,0) {
  \parbox[t][3.25cm][t]{5.65cm}{\raggedright
  \textbf{Finding Ledger}\par\smallskip
  Round 1: record \texttt{pit-001} and link to the affected output.\par\smallskip
  Round 2 omits it: the finding stays open.\par\smallskip
  Close only after the Critic verifies the repaired output.
  }
};
\end{tikzpicture}
}\end{center}
\begin{itemize}
\item Question: What check on the saved files would show that the timing problem is fixed?
\item Include every unresolved finding in the next request.
\end{itemize}
\end{frame}
\begin{frame}{Acceptance}
\begin{center}\resizebox{.95\textwidth}{!}{\begin{tikzpicture}[
  box/.style={draw=blue,fill=blue!5,rounded corners=2pt,text width=2.7cm,minimum height=1.6cm,align=center,font=\small,inner sep=5pt},
  flow/.style={-{Stealth[length=5pt]},thick,draw=black!60}]
\node[box] (g) at (1.5,0) {\textbf{Numerical gates}\\all required checks pass};
\node[box] (c) at (5,0) {\textbf{Full review}\\all critics accept; no open findings};
\node[box] (r) at (8.5,0) {\textbf{Same candidate}\\two clean rounds};
\node[box,draw=forestgreen,fill=forestgreen!6] (a) at (12,0) {\textbf{ACCEPTED}\\research task ends};
\draw[flow] (g)--(c);\draw[flow] (c)--(r);\draw[flow] (r)--(a);
\node[font=\small,text=black!60] at (6.75,-1.65) {A changed candidate starts a new clean-review streak.};
\end{tikzpicture}
}\end{center}
\end{frame}
\begin{frame}{Stopping and Recovery}
\begin{center}\resizebox{.90\textwidth}{!}{\begin{tikzpicture}[
  box/.style={draw=blue,fill=blue!5,rounded corners=3pt,text width=3.0cm,minimum height=1.05cm,align=center,font=\small,inner sep=5pt},
  good/.style={draw=forestgreen,fill=forestgreen!6},
  bad/.style={draw=darkred,fill=red!4},
  flow/.style={-{Stealth[length=5pt]},thick,draw=black!60}]
\node[box] (p) at (1.8,0) {\textbf{Observe}\\timeout or\\invalid artifact};
\node[box] (i) at (6.6,0) {\textbf{Inspect}\\completion, IDs, hashes};
\node[box,good] (r) at (11.4,0) {\textbf{Resume}\\pending work only};
\draw[flow] (p)--(i);
\draw[flow] (i)--node[above,font=\footnotesize]{verified}(r);
\node[box,bad] (f) at (6.6,-2.0) {\textbf{Stop or repair}\\blocked, changed or unsafe};
\draw[flow,draw=darkred] (i.south)--(f.north);
\end{tikzpicture}
}\end{center}
\begin{itemize}
\item Set limits on rounds, attempts, tool calls and elapsed time before starting.
\item Continue is a nonterminal route; rejection, blocking, exhaustion, stalling and user stop are distinct outcomes.
\item Check completion, repair errors, investigate changed evidence and reconcile uncertain external actions before retrying.
\item Resume only verified pending work; a completed case returns its saved result.
\end{itemize}
\end{frame}
\subsection{Graph}
\begin{frame}{Graph Roadmap}
\begin{itemize}
\item \textbf{Graph}
\begin{itemize}
\item Dependencies and handoffs
\item Artifacts and shared state
\item Teams, fan-out and joins
\item Conditional routing and combined limits
\end{itemize}
\end{itemize}
\end{frame}
\begin{frame}{Dependencies and Handoffs}
\begin{center}\resizebox{.94\textwidth}{!}{\begin{tikzpicture}[
  box/.style={draw=blue,fill=blue!5,rounded corners=2pt,text width=3.3cm,minimum height=1.5cm,align=center,font=\small,inner sep=5pt},
  flow/.style={-{Stealth[length=5pt]},thick,draw=black!60},
  lab/.style={font=\scriptsize,text=black!60,fill=white,inner sep=2pt}]
\node[box] (v) at (1.7,0) {\textbf{Validated inputs}\\dated snapshot};
\node[box] (g) at (6.6,0) {\textbf{Run candidate}\\saved outputs};
\node[box,draw=forestgreen,fill=forestgreen!6] (c) at (11.5,0) {\textbf{Review evidence}\\accept or repair};
\draw[flow] (v)--node[above,lab]{ready}(g);
\draw[flow] (g)--node[above,lab]{artifact}(c);
\draw[flow,draw=darkred] (c.south)--(11.5,-1.7)--node[below,lab,text=darkred]{failed review: return the finding}(6.6,-1.7)--(g.south);
\node[font=\small,text=black!60] at (6.6,-2.8) {Missing prerequisites: wait before dispatch.};
\end{tikzpicture}
}\end{center}
\begin{itemize}
\item Nodes perform tasks; edges specify prerequisites and handoffs.
\item Wait for missing inputs. Return failed reviews to the Generator.
\item A forward chain can be acyclic; its repair edge makes the complete workflow cyclic.
\end{itemize}
\end{frame}
\begin{frame}{Artifacts and Shared State}
\begin{center}\resizebox{.96\textwidth}{!}{\begin{tikzpicture}[
  box/.style={draw=blue,fill=blue!5,rounded corners=2pt,align=center,font=\small,inner sep=6pt},
  flow/.style={-{Stealth[length=5pt]},thick,draw=black!60}]
\node[box,text width=4.2cm,minimum height=3.1cm] (w) at (2.3,0) {\textbf{Worker artifact}\\[4pt]
request ID\\
candidate + version\\
status + findings\\
evidence path + hash};
\node[box,text width=3.3cm,minimum height=3.1cm,draw=forestgreen,fill=forestgreen!6] (v) at (7.0,0) {\textbf{Handoff validation}\\[4pt]
schema + identity\\
complete reports\\
evidence integrity};
\node[box,text width=3.5cm,minimum height=3.1cm] (s) at (11.4,0) {\textbf{Controller state}\\[4pt]
pending role + round\\
carried findings\\
next eligible task};
\draw[flow] (w)--(v);
\draw[flow] (v)--node[above,font=\scriptsize]{pass}(s);
\draw[flow,draw=darkred] (v.south)--(7.0,-2.3)--node[below,font=\scriptsize]{invalid: repair report; do not advance}(2.3,-2.3)--(w.south);
\end{tikzpicture}
}\end{center}
\begin{itemize}
\item Workers write separate reports. Only the controller updates the run's control state.
\item Validate identity, completeness and evidence before merging.
\item Keep conflicting findings explicit; the last report must not silently overwrite them.
\end{itemize}
\end{frame}
\begin{frame}{Teams and Joins}
\begin{center}\resizebox{.86\textwidth}{!}{\begin{tikzpicture}[
  lead/.style={draw,rounded corners=2pt,line width=.7pt,align=center,font=\small,inner sep=4pt},
  worker/.style={draw,dashed,rounded corners=2pt,fill=white,text width=1.30cm,minimum height=.65cm,align=center,font=\footnotesize,inner sep=2pt},
  flow/.style={-{Stealth[length=4pt]},line width=.65pt,draw=black!60},
  label/.style={font=\footnotesize,align=center,text=black!60,fill=white,inner sep=1pt}]
\node[lead,draw=forestgreen,fill=forestgreen!6,text width=5.0cm] (o) at (7.65,0)
  {\textbf{Orchestrator}\\state + gates + routing};
\node[font=\footnotesize,align=center] (idea) at (1.3,0) {\textbf{User idea}};
\draw[flow] (idea.east)--(o.west);

\draw[blue,line width=.7pt,rounded corners=3pt,fill=blue!4] (0,-1.1) rectangle (5.4,-3.95);
\draw[darkred,line width=.7pt,rounded corners=3pt,fill=red!4] (7.8,-1.1) rectangle (15.3,-3.95);
\node[font=\small\bfseries,text=blue] at (2.7,-1.38) {Generator lead};
\node[font=\small\bfseries,text=darkred] at (11.55,-1.38) {Critic lead};
\node[font=\footnotesize,text=black!60] at (2.7,-1.72) {coordinates build specialists};
\node[font=\footnotesize,text=black!60] at (11.55,-1.72) {dispatches fresh reviewers each round};

\node[worker,draw=blue] (gd) at (1.0,-2.48) {Data};
\node[worker,draw=blue] (gf) at (2.7,-2.48) {Features};
\node[worker,draw=blue] (gb) at (4.4,-2.48) {Backtest};
\node[worker,draw=darkred,text width=1.5cm,minimum height=.80cm] (cd) at (8.8,-2.55) {Data \&\\Evidence};
\node[worker,draw=darkred,text width=1.5cm,minimum height=.80cm] (cr) at (10.65,-2.55) {Robustness};
\node[worker,draw=darkred,text width=1.5cm,minimum height=.80cm] (ce) at (12.5,-2.55) {Execution};
\node[worker,draw=darkred,text width=1.5cm,minimum height=.80cm] (ck) at (14.35,-2.55) {Risk \&\\Controls};
\node[lead,draw=blue,fill=white,text width=4.3cm,minimum height=.5cm] (gi) at (2.7,-3.5)
  {\textbf{Integrate + test}};
\node[lead,draw=darkred,fill=white,text width=6.4cm,minimum height=.5cm] (cm) at (11.55,-3.5)
  {\textbf{Merge + coverage check}};
\foreach \worker in {gd,gf,gb}{\draw[flow] (2.7,-1.96)--(\worker.north);\draw[flow] (\worker.south)--(gi.north);}
\foreach \worker in {cd,cr,ce,ck}{\draw[flow] (11.55,-1.96)--(\worker.north);\draw[flow] (\worker.south)--(cm.north);}

\draw[flow] (o.south west)--(4.7,-.75)--(2.7,-.75)--(2.7,-1.1);
\draw[flow] (o.south east)--(10.3,-.75)--(11.55,-.75)--(11.55,-1.1);
\node[label] at (2.7,-.55) {dispatch};
\node[label] at (11.55,-.55) {dispatch};
\draw[flow,draw=blue] (5.4,-2.0)--(7.8,-2.0);
\node[label,text=blue] at (6.6,-1.62) {checked inputs\\via Orchestrator};
\draw[flow,dashed,draw=darkred] (7.8,-2.95)--(5.4,-2.95);
\node[circle,draw=darkred,fill=white,inner sep=1pt,text=darkred,font=\small] at (6.6,-2.95) {$\times$};
\node[label,text=darkred] at (6.6,-3.37) {reviews via\\Orchestrator};

\node[lead,draw=forestgreen,fill=forestgreen!6,text width=5.2cm] (packet) at (7.65,-4.75)
  {\textbf{Orchestrator: repair packet}\\ledger + decision: repeat / end};
\draw[flow] (cm.south)--(11.55,-4.75)--(packet.east);
\node[label,anchor=west] at (11.7,-4.35) {merged findings};
\draw[flow] (packet.west)--(-.65,-4.75)--(-.65,-1.38)--(0,-1.38);
\node[label] at (2.7,-4.35) {repair packet if continuing};
\node[font=\footnotesize,align=left,text=black!60,anchor=west] at (10.4,-5.2) {solid: role lead\\dashed: specialist};
\end{tikzpicture}
}\end{center}
\begin{itemize}
\item Use Data \& Evidence, Robustness, Execution and Risk \& Controls when checks need different expertise.
\item Wait for required reports; check identity, evidence and combined exposure before joining results.
\item Data review covers quality, provenance and timing; risk review covers limits and authorization evidence.
\end{itemize}
\end{frame}
\begin{frame}{Conditional Routing}
\begin{center}\resizebox{.82\textwidth}{!}{\begin{tikzpicture}[
  box/.style={draw=blue,fill=blue!5,rounded corners=3pt,text width=3.0cm,minimum height=1.0cm,align=center,font=\small,inner sep=5pt},
  good/.style={draw=forestgreen,fill=forestgreen!6}, bad/.style={draw=darkred,fill=red!4},
  flow/.style={-{Stealth[length=5pt]},thick,draw=black!60}]
\node[box] (w) at (1.6,0) {\textbf{Wait}\\inputs missing};
\node[box] (r) at (5.0,0) {\textbf{Repair}\\submission invalid};
\node[box,good] (n) at (8.4,0) {\textbf{Advance}\\checks pass};
\node[box,bad] (t) at (11.8,0) {\textbf{Stop}\\terminal outcome};
\draw[flow] (w)--(r);\draw[flow] (r)--(n);\draw[flow] (n)--(t);
\node[font=\scriptsize,align=center,text=black!60] at (6.7,-1.45) {A changed input version starts a new case and rebuilds dependent outputs.};
\end{tikzpicture}
}\end{center}
\begin{itemize}
\item Choose the next task using validated status and dependencies.
\item Wait when inputs are missing; repair invalid submissions at the same role; keep downstream work blocked until checks pass.
\item A changed input version starts a new case and rebuilds dependent outputs.
\item Local passes do not guarantee a valid combined exposure; the classroom fixture lacks this portfolio-wide gate.
\end{itemize}
\end{frame}
\subsection{Evolution}
\begin{frame}{Evolution Roadmap}
\begin{itemize}
\item \textbf{Evolution}
\begin{itemize}
\item What can change: prompts/context, skills, tools/code, roles/workflow
\item Evidence and evaluation
\item Regression review
\item Version control, approval and monitoring
\end{itemize}
\end{itemize}
\end{frame}
\begin{frame}{Evolution}
\begin{itemize}
\item Tool loops choose actions; research loops revise candidates.
\item Evolution evaluates reusable system changes for adoption, rejection or deferral.
\item Strategy-parameter search stays within research; QT demonstrates prompt-version evolution.
\end{itemize}
\end{frame}
\begin{frame}{Evolution Workflow}
\setlength{\topsep}{0pt}
\begin{center}
\resizebox{.90\textwidth}{!}{\begin{tikzpicture}[
  box/.style={draw=blue,fill=blue!5,rounded corners=2pt,text width=2.65cm,minimum height=1.4cm,align=center,font=\footnotesize,inner sep=4pt},
  flow/.style={-{Stealth[length=4pt]},line width=.8pt,draw=black!60}]
\node[draw=forestgreen,fill=forestgreen!5,rounded corners=2pt,font=\small,align=center,text width=12.9cm,inner sep=5pt] at (6.75,1.2)
{\textbf{Evolution Orchestrator:} failure, change scope, parent version, budget and state};
\node[box] (e) at (1.5,0) {\textbf{Completed cases}\\retained evidence};
\node[box] (d) at (5.0,0) {\textbf{Diagnose}\\recurring failure};
\node[box] (g) at (8.5,0) {\textbf{Generator team}\\scoped patch + tests};
\node[box] (t) at (12.0,0) {\textbf{Evaluate}\\fixed cases + regressions};
\node[box,draw=darkred,fill=red!4] (c) at (12.0,-1.85) {\textbf{Fresh Critic team}\\evidence + scope\\+ regressions};
\node[box,draw=darkred,fill=red!4] (h) at (8.5,-1.85) {\textbf{Owner decision}\\approve / reject\\or defer};
\node[box,draw=forestgreen,fill=forestgreen!6] (v) at (5.0,-1.85) {\textbf{Version registry}\\adopt or retain prior};
\node[box,draw=forestgreen,fill=forestgreen!6] (n) at (1.5,-1.85) {\textbf{Future cases}\\active version};
\draw[flow] (e)--(d);\draw[flow] (d)--(g);\draw[flow] (g)--(t);
\draw[flow] (t)--node[right,font=\scriptsize]{pass}(c);
\draw[flow] (c)--node[above,font=\scriptsize]{pass}(h);
\draw[flow] (h)--(v);\draw[flow] (v)--(n);
\draw[flow,dashed] (n.west)--(-.35,-1.85)--(-.35,0)--(e.west);
\node[font=\scriptsize,text=darkred,align=center] at (6.75,-2.9) {Failed checks or review: retain the active version and evidence. Rollback remains available.};
\end{tikzpicture}
}
\end{center}
\begin{itemize}
\item The Generator proposes a scoped change and tests; fresh Critics review evidence, regressions and scope.
\item Only versions passing the required checks and review can be approved. Rejected or deferred changes retain the active version.
\end{itemize}
\end{frame}
\begin{frame}{Evidence and Evaluation}
\begin{itemize}
\item Set criteria before comparing changes.
\item Check format, meaning, evidence support and authority.
\item Use development cases for feedback and held-out cases for final assessment.
\item QT timing requires execution after actual signal availability; a date field alone is insufficient.
\end{itemize}
\begin{center}\resizebox{.90\textwidth}{!}{\begin{tikzpicture}[
 box/.style={draw=blue,fill=blue!5,rounded corners=2pt,text width=5.6cm,minimum height=1.0cm,align=center,font=\small,inner sep=5pt}]
\node[box] at (3.0,0) {\textbf{Development cases}\\feedback can inform revisions};
\node[box,draw=forestgreen,fill=forestgreen!6] at (10.0,0) {\textbf{Held-out cases}\\separate final assessment};
\end{tikzpicture}
}\end{center}
\end{frame}
\begin{frame}{Regression Review}
\begin{itemize}
\item Controller tests check missing reports, stale IDs, invalid evidence and permissions.
\item Model-behavior comparisons check timing, evidence and whether the recorded failure is fixed.
\item QT's prompt interface rejects risk-limit, strategy-parameter and controller-code edits.
\item A proposal cannot weaken its own acceptance standard.
\end{itemize}
\end{frame}
\begin{frame}{Versions and Rollback}
\begin{center}
\resizebox{.94\textwidth}{!}{\begin{tikzpicture}[
 box/.style={draw=blue,fill=blue!5,rounded corners=2pt,text width=3.4cm,minimum height=1.15cm,align=center,font=\small,inner sep=5pt},
 flow/.style={-{Stealth[length=4pt]},line width=.8pt,draw=black!60},
 lab/.style={font=\scriptsize,text=black!60,fill=white,inner sep=2pt}]
\node[box] (old) at (1.8,0) {\textbf{Current version}\\retained configuration};
\node[box] (cand) at (6.6,0) {\textbf{Candidate version}\\diff + test evidence};
\node[box,draw=forestgreen,fill=forestgreen!6] (active) at (11.4,0) {\textbf{Approved version}\\new active pointer};
\draw[flow] (old)--node[above,lab]{evaluate}(cand);
\draw[flow] (cand)--node[above,lab]{approval}(active);
\node[box,fill=white,minimum height=.85cm] (existing) at (1.8,-1.8) {\textbf{Existing cases}\\keep their\\original version};
\node[box,fill=white,minimum height=.85cm] (new) at (11.4,-1.8) {\textbf{New cases}\\use the active version};
\draw[flow] (old)--(existing);\draw[flow] (active)--(new);
\draw[flow,draw=darkred,dashed] (active.south west)--(9.2,-2.95)--node[below,lab,text=darkred]{authorized rollback: restore prior pointer}(4.0,-2.95)--(old.south east);
\end{tikzpicture}
}
\end{center}
\begin{itemize}
\item Record the model, tools, data, policy and prompts tested with each version.
\item An alias selects a version; approval needs its own record.
\item Retain both versions and evidence. Rollback restores the prior pointer; verify behavior in a new case.
\end{itemize}
\end{frame}
\begin{frame}{Monitoring}
\begin{center}
\resizebox{.94\textwidth}{!}{\begin{tikzpicture}[
  box/.style={draw=blue,fill=blue!5,rounded corners=2pt,text width=3.3cm,minimum height=1.5cm,align=center,font=\small,inner sep=5pt},
  link/.style={-{Stealth[length=5pt]},thick,draw=black!60}]
\node[box] (n0) at (1.7,0) {Later cases\\new evidence};
\node[box] (n1) at (6.6,0) {Monitor failures\\and shared risks};
\draw[link] (n0) -- (n1);
\node[box,draw=forestgreen,fill=forestgreen!6] (n2) at (11.5,0) {Review changes\\before adoption};
\draw[link] (n1) -- (n2);
\end{tikzpicture}
}
\end{center}
\begin{itemize}
\item Monitor recurring failures and shared risks.
\item Re-evaluate when models, tools, data feeds or permissions change.
\item Feed recorded failures back into evolution. Prompt optimization does not guarantee reliability.
\end{itemize}
\end{frame}
\begin{frame}[t]{Complete Architecture}
\begin{center}\resizebox{!}{6.0cm}{\begin{tikzpicture}[
  lead/.style={draw,rounded corners=2pt,line width=.7pt,align=center,font=\footnotesize,inner sep=4pt},
  worker/.style={draw,dashed,rounded corners=2pt,fill=white,text width=1.30cm,minimum height=.62cm,align=center,font=\scriptsize,inner sep=2pt},
  flow/.style={-{Stealth[length=4pt]},line width=.65pt,draw=black!60},
  edge/.style={font=\scriptsize,align=center,text=black!65,fill=white,inner sep=2pt}]
\draw[darkorange!85!black,line width=1.1pt,rounded corners=5pt]
  (-.45,.88) rectangle (13.65,-5.65);
\node[anchor=west,fill=white,inner sep=4pt,font=\small\bfseries,text=darkorange!85!black]
  at (0,.88) {Evolution: evaluated system changes};
\node[lead,draw=forestgreen,fill=forestgreen!6,text width=4.5cm,minimum height=.68cm] (o) at (6.6,0)
  {\textbf{Orchestrator}\\dispatch + state + routing};
\node[font=\footnotesize,align=center] (idea) at (1.4,0) {\textbf{User idea}\\mandate + limits};
\draw[flow] (idea.east)--(o.west);
\node[font=\scriptsize,align=center,text=darkred] at (11.5,0)
  {Human policy ownership\\and approval};

\draw[blue,line width=.7pt,rounded corners=3pt,fill=blue!4] (0,-1.0) rectangle (5.4,-3.7);
\draw[darkred,line width=.7pt,rounded corners=3pt,fill=red!4] (7.8,-1.0) rectangle (13.2,-3.7);
\node[font=\footnotesize\bfseries,text=blue] at (2.7,-1.28) {Generator lead};
\node[font=\footnotesize\bfseries,text=darkred] at (10.5,-1.28) {Critic lead};
\node[font=\scriptsize,text=black!65] at (2.7,-1.60) {coordinates build specialists};
\node[font=\scriptsize,text=black!65] at (10.5,-1.60) {fresh review each round};
\node[worker,draw=blue] (gd) at (1.0,-2.25) {Data};
\node[worker,draw=blue] (gf) at (2.7,-2.25) {Features};
\node[worker,draw=blue] (gb) at (4.4,-2.25) {Backtest};
\node[worker,draw=darkred] (cp) at (8.8,-2.25) {PIT};
\node[worker,draw=darkred] (cr) at (10.5,-2.25) {Robustness};
\node[worker,draw=darkred] (ce) at (12.2,-2.25) {Execution};
\node[lead,draw=blue,fill=white,text width=4.3cm,minimum height=.52cm] (gi) at (2.7,-3.18)
  {\textbf{Integrate + test}};
\node[lead,draw=darkred,fill=white,text width=4.3cm,minimum height=.52cm] (cm) at (10.5,-3.18)
  {\textbf{Merge + coverage check}};
\foreach \w in {gd,gf,gb}{\draw[flow] (2.7,-1.83)--(\w.north);\draw[flow] (\w.south)--(gi.north);}
\foreach \w in {cp,cr,ce}{\draw[flow] (10.5,-1.83)--(\w.north);\draw[flow] (\w.south)--(cm.north);}
\draw[flow] (o.south west)--(4.0,-.68)--(2.7,-.68)--(2.7,-1.0);
\draw[flow] (o.south east)--(9.2,-.68)--(10.5,-.68)--(10.5,-1.0);
\node[edge] at (2.7,-.56) {dispatch};
\node[edge] at (10.5,-.56) {after checks};
\draw[flow,draw=blue] (5.4,-2.0)--(7.8,-2.0);
\node[edge,text=blue] at (6.6,-1.62) {executable gates\\via Orchestrator};
\node[edge,text=darkred] at (6.6,-2.85) {findings return\\to Orchestrator};
\node[lead,draw=forestgreen,fill=forestgreen!6,text width=4.6cm,minimum height=.65cm] (packet) at (6.6,-4.30)
  {\textbf{Orchestrator: repair packet}\\findings + decision: repeat / end};
\draw[flow] (cm.south)--(10.5,-4.30)--(packet.east);
\node[edge] at (10.5,-3.93) {merged findings};
\draw[flow] (packet.west)--(2.7,-4.30)--(gi.south);
\node[edge] at (2.7,-3.93) {repair if continuing};
\node[draw=black!35,rounded corners=2pt,fill=gray!5,text width=12.3cm,align=center,
  minimum height=.45cm,font=\scriptsize,inner sep=4pt] at (6.6,-5.15)
  {\textbf{Shared state and evidence:} request IDs, artifacts, hashes, findings and versions};
\end{tikzpicture}
}\end{center}
\begin{itemize}
\item Locate the role harnesses, revision loops, handoffs and evolution boundary.
\end{itemize}
\end{frame}
